% =====================================================================
% 书末：Lecture Notes in Mathematics 丛书目录（Bisher erschienen）
% 与丛书说明（英/德/法），与原书封底一致。
% =====================================================================
\cleardoublepage
\chapter*{Lecture Notes in Mathematics}
\addcontentsline{toc}{chapter}{Lecture Notes in Mathematics}
\noindent Bisher erschienen / Already published

\medskip
\begin{description}
\setlength{\itemsep}{0.15em}
\item[Vol. 1:] J. Wermer, Seminar \"uber Funktionen-Algebren. IV, 30 Seiten. 1964. DM 3,80
\item[Vol. 2:] A. Borel, Cohomologie des espaces localement compacts d'apr\`es J. Leray. IV, 93 pages. 1964. DM 9,--
\item[Vol. 3:] J. F. Adams, Stable Homotopy Theory. 2nd revised edition. IV, 78 pages. 1966. DM 7,80
\item[Vol. 4:] M. Arkowitz and C. R. Curjel, Groups of Homotopy Classes. IV, 36 pages. 1964. DM 4,80
\item[Vol. 5:] J.-P. Serre, Cohomologie Galoisienne. Troisi\`eme \'edition. VIII, 214 pages. 1965. DM 18,--
\item[Vol. 6:] H. Hermes, Eine Termlogik mit Auswahloperator. IV, 42 Seiten. 1965. DM 5,80
\item[Vol. 7:] Ph. Tondeur, Introduction to Lie Groups and Transformation Groups. VIII, 176 pages. 1965. DM 13,50
\item[Vol. 8:] G. Fichera, Linear Elliptic Differential Systems and Eigenvalue Problems. IV, 176 pages. 1965. DM 13.50
\item[Vol. 9:] P. L. Iv\u{a}nescu, Pseudo-Boolean Programming and Applications. IV, 50 pages. 1965. DM 4,80
\item[Vol. 10:] H. L\"uneburg, Die Suzukigruppen und ihre Geometrien. VI, 111 Seiten. 1965. DM 8,--
\item[Vol. 11:] J.-P. Serre, Alg\`ebre Locale Multiplicit\'es. R\'edig\'e par P. Gabriel. Seconde \'edition. VIII, 192 pages. 1965. DM 12,--
\item[Vol. 12:] A. Dold, Halbexakte Homotopiefunktoren. II, 157 Seiten. 1966. DM 12,--
\item[Vol. 13:] E. Thomas, Seminar on Fiber Spaces. VI, 45 pages. 1966. DM 4,80
\item[Vol. 14:] H. Werner, Vorlesung \"uber Approximationstheorie. IV, 184 Seiten und 10 Seiten Anhang. 1966. DM 14,--
\item[Vol. 15:] F. Oort, Commutative Group Schemes. VI, 133 pages. 1966. DM 9,80
\item[Vol. 16:] J. Pfanzagl and W. Pierlo, Compact Systems of Sets. IV, 48 pages. 1966. DM 5,80
\item[Vol. 17:] C. M\"uller, Spherical Harmonics. IV, 46 pages. 1966. DM 5,--
\item[Vol. 18:] H.-B. Brinkmann, Kategorien und Funktoren. Nach einer Vorlesung von D. Puppe. XII, 107 Seiten. 1966. DM 8,--
\item[Vol. 19:] G. Stolzenberg, Volumes, Limits and Extensions of Analytic Varieties. IV, 47 pages. 1966. DM 5,40
\item[Vol. 20:] R. Hartshorne, Residues and Duality. VIII, 423 pages. 1966. DM 20,--
\item[Vol. 22:] H. Bauer, Harmonische R\"aume und ihre Potentialtheorie. IV, 175 Seiten. 1966. DM 14,--
\item[Vol. 23:] P. L. Iv\u{a}nescu and S. Rudeanu, Pseudo-Boolean Methods for Bivalent Programming. VI, 113 pages. 1966. DM 10,--
\end{description}

\clearpage

% ---------- 英文说明 ----------
\thispagestyle{plain}
This series aims at speedy, informal, and high level information on new developments in mathematical research and teaching. Considered for publication are:
\begin{enumerate}[leftmargin=2.5em]
\item Preliminary drafts of original papers and monographs
\item Special lectures on a new field, or a classical field from a new point of view
\item Seminar reports
\item Reports from meetings
\end{enumerate}
Out of print manuscripts satisfying the above characterization may also be considered, if they continue to be in demand.

\medskip
The timeliness of a manuscript is more important than its form, which may be unfinished and preliminary. In certain instances, therefore, proofs may only be outlined, or results may be presented which have been or will also be published elsewhere.

\medskip
The publication of the ``Lecture Notes'' Series is intended as a service, in that a commercial publisher, Springer-Verlag, makes house publications of mathematical institutes available to mathematicians on an international scale. By advertising them in scientific journals, listing them in catalogs, further by copyrighting and by sending out review copies, an adequate documentation in scientific libraries is made possible.

\section*{Manuscripts}
Since manuscripts will be reproduced photomechanically, they must be written in clean typewriting. Handwritten formulae are to be filled in with indelible black or red ink. Any corrections should be typed on a separate sheet in the same size and spacing as the manuscript. All corresponding numerals in the text and on the correction sheet should be marked in pencil. Springer-Verlag will then take care of inserting the corrections in their proper places. Should a manuscript or parts thereof have to be retyped, an appropriate indemnification will be paid to the author upon publication of his volume. The authors receive $25$ free copies.

\medskip
Manuscripts in English, German or French should be sent to Prof.\ Dr.\ A.\ Dold, Mathematisches Institut der Universit\"at Heidelberg, Tiergartenstra\ss{}e or Prof.\ Dr.\ B.\ Eckmann, Eidgen\"ossische Technische Hochschule, Z\"urich.

% ---------- 德文说明 ----------
\section*{\texorpdfstring{Die „Lecture Notes“}{Die Lecture Notes}}
Die „Lecture Notes“ sollen rasch und informell, aber auf hohem Niveau, \"uber neue Entwicklungen der mathematischen Forschung und Lehre berichten. Zur Ver\"offentlichung kommen:
\begin{enumerate}[leftmargin=2.5em]
\item Vorl\"aufige Fassungen von Originalarbeiten und Monographien.
\item Spezielle Vorlesungen \"uber ein neues Gebiet oder ein klassisches Gebiet in neuer Betrachtungsweise.
\item Seminarausarbeitungen.
\item Vortr\"age von Tagungen.
\end{enumerate}
Ferner kommen auch \"altere vergriffene spezielle Vorlesungen, Seminare und Berichte in Frage, wenn nach ihnen eine anhaltende Nachfrage besteht.

\medskip
Die Beitr\"age d\"urfen im Interesse einer gr\"o\ss{}eren Aktualit\"at durchaus den Charakter des Unfertigen und Vorl\"aufigen haben. Sie brauchen Beweise unter Umst\"anden nur zu skizzieren und d\"urfen auch Ergebnisse enthalten, die in \"ahnlicher Form schon erschienen sind oder sp\"ater erscheinen sollen.

\medskip
Die Herausgabe der „Lecture Notes“ Serie durch den Springer-Verlag stellt eine Dienstleistung an die mathematischen Institute dar, indem der Springer-Verlag f\"ur ausreichende Lagerhaltung sorgt und einen gro\ss{}en internationalen Kreis von Interessenten erfassen kann. Durch Anzeigen in Fachzeitschriften, Aufnahme in Kataloge und durch Anmeldung zum Copyright sowie durch die Versendung von Besprechungsexemplaren wird eine l\"uckenlose Dokumentation in den wissenschaftlichen Bibliotheken erm\"oglicht.

\section*{Beschaffenheit der Manuskripte}
Die Manuskripte werden photomechanisch vervielf\"altigt; sie m\"ussen daher in sauberer Schreibmaschinenschrift geschrieben sein. Handschriftliche Formeln bitte nur mit schwarzer Tusche oder roter Tinte eintragen. Korrekturw\"unsche werden in der gleichen Maschinenschrift auf einem besonderen Blatt erbeten (Zuordnung der Korrekturen im Text und auf dem Blatt sind durch Bleistiftziffern zu kennzeichnen). Der Verlag sorgt dann f\"ur das ordnungsgem\"a\ss{}e Tektieren der Korrekturen. Falls das Manuskript oder Teile desselben neu geschrieben werden m\"ussen, ist der Verlag bereit, dem Autor bei Erscheinen seines Bandes einen angemessenen Betrag zu zahlen. Die Autoren erhalten 25 Freiexemplare.

\medskip
Manuskripte, in englischer, deutscher oder franz\"osischer Sprache abgefa\ss{}t, nimmt Prof.\ Dr.\ A.\ Dold, Mathematisches Institut der Universit\"at Heidelberg, Tiergartenstra\ss{}e oder Prof.\ Dr.\ B.\ Eckmann, Eidgen\"ossische Technische Hochschule, Z\"urich, entgegen.

% ---------- 法文说明 ----------
\section*{\texorpdfstring{Cette s\'erie}{Cette serie}}
\setlength{\emergencystretch}{3em}
Cette s\'erie a pour but de donner des informations rapides, de niveau \'elev\'e, sur des d\'eveloppements r\'ecents en math\'ematiques, aussi bien dans la recherche que dans l'enseignement sup\'erieur. On pr\'evoit de publier
\begin{enumerate}[leftmargin=2.5em]
\item des versions pr\'eliminaires de travaux originaux et de monographies
\item des cours sp\'eciaux portant sur un domaine nouveau ou sur des aspects nouveaux de domaines classiques
\item des rapports de s\'eminaires
\item des conf\'erences faites \`a des congr\`es ou des colloquiums
\end{enumerate}
En outre il est pr\'evu de publier dans cette s\'erie, si la demande le justifie, des rapports de s\'eminaires et des cours multicopi\'es ailleurs qui sont \'epuis\'es.

\medskip
Dans l'int\'er\^et d'une grande actualit\'e les contributions pourront souvent \^etre d'un caract\`ere provisoire; le cas \'ech\'eant, les d\'emonstrations ne seront donn\'ees qu'en grande ligne, et les r\'esultats et m\'ethodes pourront \'egalement para\^itre ailleurs. Par cette s\'erie de »pr\'epublications« les \'editeurs Springer esp\`erent rendre d'appr\'eciables services aux instituts de math\'ematiques par le fait qu'une r\'eserve suffisante d'exemplaires sera toujours \`a disposition et que les int\'eress\'es pourront plus facilement \^etre atteints. Les annonces dans les revues sp\'ecialis\'ees, les inscriptions aux catalogues et les copyrights faciliteront pour les biblioth\`eques math\'ematiques la t\^ache de dresser une documentation compl\`ete.

\section*{Pr\'esentation des manuscrits}
Les manuscrits \'etant reproduits par proc\'ed\'e photom\'ecanique, doivent \^etre soigneusement dactylographi\'es. Il est demand\'e d'\'ecrire \`a l'encre de Chine ou \`a l'encre rouge les formules non dactylographi\'ees. Des corrections peuvent \'egalement \^etre dactylographi\'ees sur une feuille s\'epar\'ee (pri\`ere d'indiquer au crayon leur ordre de classement dans le texte et sur la feuille), la maison d'\'edition se chargeant ensuite de les ins\'erer \`a leur place dans le texte. S'il s'av\`ere n\'ecessaire d'\'ecrire de nouveau le manuscrit, soit compl\`etement, soit en partie, la maison d'\'edition se d\'eclare pr\^ete \`a se charger des frais \`a la parution du volume. Les auteurs recoivent 25 exemplaires gratuits.

\medskip
Les manuscrits en anglais, allemand ou fran\c cais peuvent \^etre adress\'es au Prof.\ Dr.\ A.\ Dold, Mathematisches Institut der Universit\"at Heidelberg, Tiergartenstra\ss{}e ou Prof.\ Dr.\ B.\ Eckmann, Eidgen\"ossische Technische Hochschule, Z\"urich.
